
\subsection{Tarski's well ordering theorem}
Assuming that for every set $A$, assuming it's infinite, then $A \times A \sim A$, we can prove that the axiom of choice holds true by first 
proving that every set can be well ordered. We will hence discuss the Tarski's proof of the well ordering theorem.

\begin{theorem}[Tarski's well ordering theorem]
    Assuming that every infinite set $A$ has the property that $A \times A \sim A$, then every set can be well ordered.
\end{theorem}
\begin{proof}
    Of course every finite set can be well ordered, since a bijective function between such set and a proper finite subset of 
    the natural numbers can be created, and since every finite subset of natural numbers is also well ordered. Hence, whithout 
    loss of generality, we can assume that the set $A$ is infinite. \\

    Let $A$ be an infinite set and let $\beta$ be the Hartogs ordinal for $A$. Wich means that $\beta$ is the smallest ordinal
    that cannot be injected into $A$ (So the smallest ordinal that is greater then $|A|$). Such ordinal can be constructed 
    as the transitive set of all ordinals with cardinality greater then or equal to $|A|$ as shown in the proof of 
    the Hartogs theorem. \\

    Let $A' := A \cup \beta$ be the union of the infinite set $A$ and its Hartogs ordinal. Since $A$ is infinite, then $A'$ is
    also infinite, and since $A'$ is infinite, under our assumption, $A' \times A' \sim A'$. \\

    For every $a \in A$ we can construct the family $S_a$ which has some intresting properties. \\

    $$ S_a := \vv{f}(\{a\} \times \beta) \cap \beta $$ \\

    Since $f$ is a bijection, then the images of two disjoint sets are also disjoint. This 
    means that the family $S_a$ is a family of disjoint 
    sets. Since $|(\{a\} \times \beta)| = |\beta| > |A|$ then its image trough $f$ is also 
    greater then $|A|$ given that $f$ is a bijection. \\

    $$ |\vv{f}(\{a\} \times \beta)| > |A| $$
    $$ \vv{f}(\{a\} \times \beta) \subseteq A \cup \beta $$
    $$ \therefore \vv{f}(\{a\} \times \beta) \cap \beta \neq \varnothing $$ \\

    We can finally define the relation 
    $a \leq b \iff \texttt{min}(S_a) \leq \texttt{min}(S_b)$ which 
    is a total ordering over the set $A$. We can then assert that since the set $A$ is 
    well-ordered this means that $A$ is in bijection with some ordinal.  \\

\end{proof}